\documentclass[12pt,a4paper,twoside]{report}

% Wybierz styl nagłówka rozdziału:
% 'classic' — Rozdział X (nad tytułem)
% 'inline'  — X. Tytuł rozdziału (w jednej linii)

\def\chapterstyletype{classic}  % lub zmień na 'inline'

%MARGINESY
\usepackage[left=3cm, right=2.5cm, top=2.5cm, bottom=2.5cm, headsep=1.2cm]{geometry}

%INTERLINIA
\linespread{1.3}

%Język polski
\usepackage{polski}
\usepackage[utf8]{inputenc}
\usepackage[T1]{fontenc}

% Pakiet, który poprawia typografię: kerning, justowanie i drobne szczegóły tekstu
\usepackage{microtype}

%landscape obrocone
\usepackage{pdflscape}

%tekst matematyczny
\usepackage{amsmath}
\usepackage{amsfonts}
\usepackage{amssymb}
\usepackage{amsthm}

%dodawanie grafiki
\usepackage{graphicx}

%do grupowania rysunków
\usepackage{subcaption}

%brak wdów - pojedyncze linie
\usepackage[all]{nowidow}

%Zapobieganie dzieleniu się wyrazów
%\tolerance=1
\emergencystretch=\maxdimen
\hyphenpenalty=10000
\hbadness=10000

%do kolorów
%\usepackage{color}
\usepackage{xcolor}

\usepackage{textcomp}%do symbolu euro
\usepackage[official]{eurosym}%drugi symbol euro

\usepackage[colorlinks=true,allcolors=black]{hyperref}

\usepackage{indentfirst} %dodaje wcięcie w pierwszym akapicie

%listy numerowane
\usepackage{enumitem}
\setlist[enumerate]{noitemsep,nolistsep}
\setlist[itemize]{noitemsep,nolistsep} 

%do modyfikacji wyglądu spisu treści, listy rysunków i tabel
\usepackage{tocloft}

%do tabel
\usepackage{tabularx}
\usepackage{multirow} %do połączenia wierszy
\usepackage{diagbox}%do ukośnej linii w tabelii
\usepackage{booktabs}
\usepackage{array}
\newcolumntype{Y}{>{\centering\arraybackslash}X}
\newcolumntype{j}{>{\raggedleft}p{200pt}}
\usepackage{caption}
\usepackage{lscape}





%podpisy rysunków i tabel
\renewcommand*{\figurename}{Rys.} %zmiana podpisu rysunku
\renewcommand*{\tablename}{Tab.} %zmiana podpisu tabeli
\captionsetup{font=small, labelsep=period, justification=centerlast} % Modyfikacja stylu podpisów: 
% - font=small ustawia rozmiar czcionki podpisu na 11 pt (można użyć footnotesize dla 10 pt),
% - labelsep=period wprowadza kropkę po numerze (zamiast dwukropka),
% - justification=centerlast centrowanie podpisu z wyrównaniem ostatniego wiersza do środka.


%%%%%%%%%%%%%%%%%%%%%%%%%%%%%%%%%%%%%%%%%%%%%%%%%%%%%%%%%%
% Ustawienie niestandardowej numeracji sekcji i podsekcji.
% Dodanie kropki po numerach (np. z "2.1 Rozdział" na "2.1. Rozdział").
% Można odkomentować, jeśli chcemy standardową formę numeracji - wtedy odkomentować wszystko poniżej do %%spis treści%%%%
\renewcommand{\thesection}{\arabic{chapter}.\arabic{section}.} % Numeracja sekcji w formacie: numer_rozdziału.numer_sekcji (np. 2.3.)
\renewcommand{\thesubsection}{\arabic{chapter}.\arabic{section}.\arabic{subsection}.} % Numeracja podsekcji w formacie: numer_rozdziału.numer_sekcji.numer_podsekcji (np. 2.3.1.)




%%%%%%%%%%%%%%%%%%%%%%%%Spis treści%%%%%%%%%%%%%%%%%%%%%%%%%%%%%
\setlength{\cftbeforechapskip}{0pt} %ustawia odstęp przed rozdziałami w spisie treści na 0pt; można zmienić wartość (np. 10pt), by zwiększyć odstęp, lub zakomentować tę linię, by przywrócić domyślne ustawienie

%%%%%%%%%%%%%%%%%%%%%%%%Tytuły rozdziałów%%%%%%%%%%%%%%%%%%%%%%%
\usepackage{titlesec}
\renewcommand{\cftchapaftersnum}{.}%ustawia, że po numerze rozdziału w spisie treści pojawi się kropka .
%\renewcommand{\chaptername}{}
%\renewcommand{\thechapter}{\arabic{chapter}.} % Dodaje kropkę po numerze rozdziału

\usepackage{ifthen}
\usepackage{ragged2e}  % do lewego wyrównania z lepszym łamaniem niż standardowy \raggedright

%instrukcja warunkowa pozwalająca na wybór formatowania nagłówka/tytułu rozdziału

\ifthenelse{\equal{\chapterstyletype}{classic}}{
  % Styl: Rozdział X (na górze), tytuł pod spodem
  \titleformat{\chapter}[display]
    {\normalfont\LARGE\bfseries\RaggedRight}
    {\chaptertitlename\ \thechapter} % np. "Rozdział 1"
    {0pt}
    {\LARGE}
}{%
  % Styl: X. Tytuł rozdziału (inline) z wiszącym wcięciem i wyrównaniem wielowierszowym
  \titleformat{\chapter}[block]
    {\normalfont\LARGE\bfseries\RaggedRight}
    {}
    {0pt}
    {%
      \makebox[1em][l]{\thechapter.}%
      \hangindent=1em
      \hangafter=1
      \ignorespaces
    }
}

% Styl rozdziału bez numeru (dla Bibliografii, Streszczenia itd.)
\titleformat{name=\chapter, numberless}[block]
  {\normalfont\LARGE\bfseries\RaggedRight}
  {}
  {0pt}
  {\hangindent=1em\hangafter=1}

%styl sekcji
\titleformat{\section}[hang]
  {\normalfont\Large\bfseries\RaggedRight} % format
  {\thesection}               % section number with dot
  {0.3em}                       % space between number and title
  {}                          % before title code
%styl podsekcji
  \titleformat{\subsection}[hang]
  {\normalfont\large\bfseries\RaggedRight} % format
  {\thesubsection}               % section number with dot
  {0.3em}                       % space between number and title
  {}                          % before title code
%stylpodpodsekcji
  \titleformat{\subsubsection}[hang]
  {\normalfont\normalsize\bfseries} % format
  {}               % brak numeru, czyli pusty label
  {0pt}            % odstęp między label a tytułem (brak, bo label pusty)
  {}                % before title code

% Ustawienie odstępów wokół tytułu rozdziału:
% \titlespacing*{<komenda>}{<lewy margines>}{<przed tytułem>}{<po tytule>}
% Tutaj:
% - lewy margines       = 0pt (brak wcięcia)
% - odstęp przed tytułem = 0pt (brak pustego miejsca na górze)
% - odstęp po tytule     = \baselineskip (dokładnie jeden wiersz odstępu przed tekstem)
\titlespacing*{\chapter}{0pt}{0pt}{\baselineskip}


%definicje, twierdzenia,...
\newtheorem{twierdzenie}{Twierdzenie}[chapter]
\newtheorem{lemat}[twierdzenie]{Lemat}
\newtheorem{wniosek}[twierdzenie]{Wniosek}
\newtheorem{obserwacja}[twierdzenie]{Obserwacja}
\newtheorem{uwaga}[twierdzenie]{Uwaga}
\theoremstyle{definition} %prostowanie czcionki
\newtheorem{przyklad}[twierdzenie]{Przykład}
\newtheorem{definicja}[twierdzenie]{Definicja}


%używane pakiety
%\usepackage{setspace}
%\usepackage{cite}
%\usepackage{tikz}
%\usepackage{subcaption}


\usepackage{algorithm}
\usepackage{algpseudocode}
\usepackage{listings}






\definecolor{niebieski}{RGB}{0,130,150} %kolor niebieski
\definecolor{mygray}{gray}{0.6} 		% kolor szary
\definecolor{zielony}{rgb}{0.98, 1.0, 0.98} % kolor zielony

\lstset{ 
	backgroundcolor=\color{zielony},   % choose the background color; you must add \usepackage{color} or \usepackage{xcolor}; should come as last argument
	basicstyle=\sffamily, %\footnotesize\ttfamily\fontsize{11}{0.8em}\selectfont,        % the size of the fonts that are used for the code
	breakatwhitespace=false,         % sets if automatic breaks should only happen at whitespace
	breaklines=true,                 % sets automatic line breaking
	captionpos=b,                    % sets the caption-position to bottom
	columns=fullflexible,
	commentstyle=\color{mygreen},    % comment style
	escapeinside={\%*}{*)},          % if you want to add LaTeX within your code
	extendedchars=true,              % lets you use non-ASCII characters; for 8-bits encodings only, does not work with UTF-8
	firstnumber=1,                   % start line enumeration with line 1
	frame=single,	                   % adds a frame around the code
	keepspaces=true,                 % keeps spaces in text, useful for keeping indentation of code (possibly needs columns=flexible)
	keywordstyle=\color{blue},       % keyword style
	numbers=left,                    % where to put the line-numbers; possible values are (none, left, right)
	numbersep=5pt,                   % how far the line-numbers are from the code
	numberstyle=\tiny\color{mygray}, % the style that is used for the line-numbers
	rulecolor=\color{black},         % if not set, the frame-color may be changed on line-breaks within not-black text (e.g. comments (green here))
	xleftmargin=0pt,%\parindent,
	showspaces=false,                % show spaces everywhere adding particular underscores; it overrides 'showstringspaces'
	showstringspaces=false,          % underline spaces within strings only
	showtabs=false,                  % show tabs within strings adding particular underscores
	language=python,
	morekeywords={*, end, function, ...}
	stepnumber=1,                    % the step between two line-numbers. If it's 1, each line will be numbered
	stringstyle=\color{mymauve},     % string literal style
	tabsize=2	                   % sets default tabsize to 2 spaces
}



%parametry dla kodów źródłowych
\definecolor{mygreen}{rgb}{0,0.6,0}
\definecolor{mygray}{rgb}{0.5,0.5,0.5}
\definecolor{mymauve}{rgb}{0.58,0,0.82}




% --- robocze: usunąć przed oddaniem pracy (razem z linią % --- robocze: usunąć przed oddaniem pracy (razem z linią % --- robocze: usunąć przed oddaniem pracy (razem z linią \input{makra} w pliku głównym) ---
% Znacznik miejsca do weryfikacji, widoczny w PDF.
\newcommand{\dower}[1]{\textcolor{red}{[DO WERYFIKACJI: #1]}}
% Plan podrozdziału (z docs/02_spis_tresci.md), widoczny w PDF szarym drukiem.
\newenvironment{planroboczy}[1][]{%
  \par\begingroup\small\color{black!60}%
  \noindent\textbf{Plan roboczy\if\relax\detokenize{#1}\relax\else\ (#1)\fi}\par\nopagebreak
}{\par\endgroup\medskip}
\newcommand{\zrodlaplanowane}[1]{\par\noindent\textbf{Źródła planowane:} #1\par}
\newcommand{\materialyrobocze}[1]{\par\noindent\textbf{Materiały robocze:} #1\par}
 w pliku głównym) ---
% Znacznik miejsca do weryfikacji, widoczny w PDF.
\newcommand{\dower}[1]{\textcolor{red}{[DO WERYFIKACJI: #1]}}
% Plan podrozdziału (z docs/02_spis_tresci.md), widoczny w PDF szarym drukiem.
\newenvironment{planroboczy}[1][]{%
  \par\begingroup\small\color{black!60}%
  \noindent\textbf{Plan roboczy\if\relax\detokenize{#1}\relax\else\ (#1)\fi}\par\nopagebreak
}{\par\endgroup\medskip}
\newcommand{\zrodlaplanowane}[1]{\par\noindent\textbf{Źródła planowane:} #1\par}
\newcommand{\materialyrobocze}[1]{\par\noindent\textbf{Materiały robocze:} #1\par}
 w pliku głównym) ---
% Znacznik miejsca do weryfikacji, widoczny w PDF.
\newcommand{\dower}[1]{\textcolor{red}{[DO WERYFIKACJI: #1]}}
% Plan podrozdziału (z docs/02_spis_tresci.md), widoczny w PDF szarym drukiem.
\newenvironment{planroboczy}[1][]{%
  \par\begingroup\small\color{black!60}%
  \noindent\textbf{Plan roboczy\if\relax\detokenize{#1}\relax\else\ (#1)\fi}\par\nopagebreak
}{\par\endgroup\medskip}
\newcommand{\zrodlaplanowane}[1]{\par\noindent\textbf{Źródła planowane:} #1\par}
\newcommand{\materialyrobocze}[1]{\par\noindent\textbf{Materiały robocze:} #1\par}
 % robocze: makra \dower i planroboczy; usunąć przed oddaniem pracy

\begin{document}

%-----------------------------------------------------------------------------------
%                            STRONA TYTUŁOWA
%-----------------------------------------------------------------------------------
% W zależności od potrzeb można załadować jedną z dwóch wersji strony tytułowej:
%
% 1. strona_tytulowa-1
%    - Klasyczna wersja strony tytułowej
%    - Tytuł pracy oraz dane autora/promotora są wyśrodkowane
%    - NIE jest dostosowana do okładki PRz (brak zgodności z okienkiem)
%
% 2. strona_tytulowa-2
%    - Wersja dopasowana do obowiązującego wzoru okładki PRz
%    - Tytuł, autor, promotor są umieszczone tak, aby były widoczne
%      w okienku okładki zgodnej z Systemem Identyfikacji Wizualnej PRz
%
% ABY WYBRAĆ WŁAŚCIWĄ STRONĘ TYTUŁOWĄ:
% - pozostaw aktywną linię \thispagestyle{empty}
\linespread{1.0}
\newgeometry{top=2.5cm, bottom=2.5cm, left=3cm, right=3cm}

	\begin{titlepage}
		\def\dystans{1cm}
		\begin{center}
			\Large 
\begin{figure}[h]
    \centering
    \includegraphics[height=3.2cm]{prz_pl.png}%
    \hfill
    \includegraphics[height=3.2cm]{wmifs_pl.png}
\end{figure} 
			{\large\scshape Politechnika Rzeszowska}
            
            {\large im. Ignacego Łukasiewicza}
        
            {\large\scshape\MakeUppercase{Wydział Matematyki i Fizyki Stosowanej}}
		\end{center}
		\vspace{1.3cm}
        
\begin{center}

        \begin{center}
            {\Large Kamil Piejko}
        \end{center}    
        \begin{center}
            {\Large\textbf{\textsc{Wpływ promptu na progi bólu agenta LLM}}}
        \end{center}
        \begin{center}
            Promotor: dr inż. Marcin Kowalik, prof. PRz
        \end{center}

\end{center}

\vspace{2cm}

        
		\begin{center}
			{\large  \textsc{Praca magisterska} \\
				\vspace{0.5cm}
				\large kierunek studiów: Inżynieria i~Analiza Danych}
		\end{center}
		\vspace{1.1cm}
		
		

		
		
		\vfill 
		\begin{center}
			{\large Rzeszów 2027}
		\end{center}
	\end{titlepage}
    \restoregeometry lub
% - zamień ją na:        \input{strona_tytulowa-2}
%
% UWAGA: Nie należy używać obu jednocześnie!

    
    \thispagestyle{empty}
\linespread{1.0}
\newgeometry{top=2.5cm, bottom=2.5cm, left=3cm, right=3cm}

	\begin{titlepage}
		\def\dystans{1cm}
		\begin{center}
			\Large 
\begin{figure}[h]
    \centering
    \includegraphics[height=3.2cm]{prz_pl.png}%
    \hfill
    \includegraphics[height=3.2cm]{wmifs_pl.png}
\end{figure} 
			{\large\scshape Politechnika Rzeszowska}
            
            {\large im. Ignacego Łukasiewicza}
        
            {\large\scshape\MakeUppercase{Wydział Matematyki i Fizyki Stosowanej}}
		\end{center}
		\vspace{1.3cm}
        
\begin{center}

        \begin{center}
            {\Large Kamil Piejko}
        \end{center}    
        \begin{center}
            {\Large\textbf{\textsc{Wpływ promptu na progi bólu agenta LLM}}}
        \end{center}
        \begin{center}
            Promotor: dr inż. Marcin Kowalik, prof. PRz
        \end{center}

\end{center}

\vspace{2cm}

        
		\begin{center}
			{\large  \textsc{Praca magisterska} \\
				\vspace{0.5cm}
				\large kierunek studiów: Inżynieria i~Analiza Danych}
		\end{center}
		\vspace{1.1cm}
		
		

		
		
		\vfill 
		\begin{center}
			{\large Rzeszów 2027}
		\end{center}
	\end{titlepage}
    \restoregeometry %lub strona_tytulowa-2
    \linespread{1.3}
    \newpage 
%----------------------------------------------------------------------------------
%                       PUSTA STRONA PO STRONIE TYTUŁOWEJ
%----------------------------------------------------------------------------------


    \thispagestyle{empty} %pusta strona bez numeru po stronie tytułowej
    ~\vfill
	\newpage
    
%---------------------------------------------------------------------------------
%                      STRONA Z PODZIĘKOWANIEM
%---------------------------------------------------------------------------------
\vspace*{16cm}
\thispagestyle{empty}
\begin{flushright}
\textit{Miejsce na ewentualne podziękowania.\\
W przypadku braku podziękowań, tę stronę wraz z następującą po niej pustą stroną\\
należy usunąć z wersji końcowej pracy.}%-- -Wypełnia student. Jeśli nie ma podziękowań, to można odkomentować stronę z podziękowaniem i stronę po stronie z podziękowaniem.
\end{flushright}
\newpage

%----------------------------------------------------------------------------------
%                       PUSTA STRONA PO STRONIE Z PODZIĘKOWANIEM
%----------------------------------------------------------------------------------

\thispagestyle{empty}
~\vfill
\newpage

%---------------------------------------------------------------------------------
%                      SPIS TREŚCI
%---------------------------------------------------------------------------------
	\renewcommand{\contentsname}{\LARGE Spis treści}
    \tableofcontents
    \setcounter{tocdepth}{1}
    %\setcounter{page}{1} %numeracja stron od 1
	%\thispagestyle{headings}
    \cleardoublepage %wymusza rozpoczęcie nowego rozdziału na stronie nieparzystej (prawej); dodaje pustą stronę, jeśli potrzeba
    %\clearpage  %kończy bieżącą stronę i wypisuje wszystkie zaległe elementy (np. rysunki, tabele), ale nie wymusza rozpoczęcia nowego rozdziału na stronie nieparzystej
    %\newpage %Wymusza przejście do nowej strony, ale nie kończy otwartych elementów

%---------------------------------------------------------------------------------
%                      WSTĘP
%---------------------------------------------------------------------------------
	
	\chapter*{Wstęp} 
    \addcontentsline{toc}{chapter}{Wstęp}
	%\thispagestyle{empty}%strona nie jest numerowana
    %% Szkielet wygenerowany z docs/02_spis_tresci.md (plan v1, 11.11.2026).
%% Blok planroboczy usuń, gdy podrozdział jest napisany.

\begin{planroboczy}[ok. 2 s.]
Problem w~trzech akapitach:
\begin{enumerate}
  \item agenci LLM podejmują decyzje w~środowiskach z~ograniczonymi zasobami,
  \item model uczenia motywowanego daje język bólu i~progu,
  \item nie wiadomo, jak bardzo prompt przesuwa te progi.
\end{enumerate}

Dalej: cel pracy, pytania badawcze (skrót z \texttt{01}), wkład własny i~krótki przegląd rozdziałów. Wymóg WMiFS: we wstępie kilka zdań o~uzyskanych wynikach, więc wstęp piszemy na końcu.

\zrodlaplanowane{\cite{starzyk2017needs, wang2024survey, sclar2024quantifying, zawislak2025budowa}}
\end{planroboczy}


    \cleardoublepage %

%%%%%%%%%%%%%%%%%%%%% --- Zasadnicza treść rozprawy zaczyna się tutaj --- %%%%%%%%%%%%%%%%%%%%%

    %% Szkielet wygenerowany z docs/02_spis_tresci.md (plan v1, 11.11.2026).
%% Blok planroboczy usuń, gdy podrozdział jest napisany.

\chapter{Motywacja i~ból w~systemach autonomicznych}
\label{ch:motywacja}

% objętość planowana: ok. 8 s.

\section{Agent, środowisko i~cel}
\label{sec:1-1}

\begin{planroboczy}[ok. 1 s.]
\begin{itemize}
  \item Definicja agenta i~środowiska, stan, akcja, polityka.
  \item Tylko tyle, ile potrzeba do dalszych rozdziałów. Bez rozbudowanego wykładu o~RL, bo nie jest przedmiotem pracy.
\end{itemize}

\zrodlaplanowane{\cite{russell2020artificial, sutton2018reinforcement, astrom2021feedback, wang2024survey}}
\end{planroboczy}

\section{Motywacja w~psychologii: potrzeby, popęd, homeostaza}
\label{sec:1-2}

\begin{planroboczy}[ok. 1,5 s.]
\begin{itemize}
  \item Potrzeby i~ich hierarchia (Maslow), motywacja wewnętrzna i~zewnętrzna (teoria autodeterminacji).
  \item Homeostaza jako utrzymywanie zmiennych wewnętrznych w~zakresie. Redukcja odchylenia jako źródło motywacji.
  \item Wniosek: ból to sygnał odchylenia od stanu pożądanego.
\end{itemize}

\zrodlaplanowane{\cite{maslow1943theory, ryan2000self, keramati2014homeostatic, yoshida2025linking}}
\end{planroboczy}

\section{Motywacja wewnętrzna w~uczeniu maszynowym}
\label{sec:1-3}

\begin{planroboczy}[ok. 1,5 s.]
\begin{itemize}
  \item Typologia podejść obliczeniowych (ciekawość, kompetencja).
  \item Motywacja wewnętrzna w~RL.
  \item Homeostatyczne RL jako most między nagrodą a~stabilnością fizjologiczną.
  \item Krótko, jako tło dla podejścia Starzyka.
\end{itemize}

\zrodlaplanowane{\cite{oudeyer2007intrinsic, singh2010intrinsically, keramati2014homeostatic, du2023guiding, klissarov2024motif, yoshida2025linking}}
\end{planroboczy}

\section{Uczenie motywowane: potrzeby, ból prymitywny i~abstrakcyjny, próg bólu}
\label{sec:1-4}

\begin{planroboczy}[ok. 3 s.]
Rdzeń teoretyczny pracy.
\begin{itemize}
  \item Agent motywowany Starzyka: potrzeby predefiniowane, ból prymitywny jako miara odległości od zaspokojenia potrzeby.
  \item Próg: agent działa tylko wtedy, gdy ból przekracza próg.
  \item Bóle abstrakcyjne tworzone w~toku uczenia (zasoby potrzebne do zaspokojenia potrzeb).
  \item Wybór celu przez rywalizację bólów (WTA).
  \item Zachowania oportunistyczne. Architektura MLECOG.
  \item Rysunek: schemat zależności ból–akcja w~środowisku ogrodowym (na podstawie grafu z~pracy referencyjnej, przerysowany).
  \item Wniosek dla pracy: w~agencie Starzyka próg jest parametrem, w~agencie LLM trzeba go wyznaczyć z~zachowania.
\end{itemize}

\zrodlaplanowane{\cite{starzyk2008motivation, starzyk2011motivated, starzyk2012motivated, starzyk2013simulation, graham2015opportunistic, starzyk2017mlecog, starzyk2017needs}}
\end{planroboczy}

\section{Sztuczny ból i~homeostaza w~robotyce}
\label{sec:1-5}

\begin{planroboczy}[ok. 1 s.]
\begin{itemize}
  \item Przykłady „bólu” w~robotach: odruchy na szkodliwy kontakt, homeostaza w~projektowaniu maszyn „czujących”.
  \item Rozróżnienie: ból funkcjonalny (sygnał sterujący) a~doznanie. Praca zajmuje się wyłącznie pierwszym.
\end{itemize}

\zrodlaplanowane{\cite{kuehn2017artificial, man2019homeostasis, sharkey2025could, asada2019artificial}}
\end{planroboczy}

    \cleardoublepage

    %% Szkielet wygenerowany z docs/02_spis_tresci.md (plan v1, 11.11.2026).
%% Blok planroboczy usuń, gdy podrozdział jest napisany.

\chapter{Duże modele językowe jako agenci decyzyjni}
\label{ch:llm-agenci}

% objętość planowana: ok. 9 s.

\section{Model językowy, generowanie i~próbkowanie}
\label{sec:2-1}

\begin{planroboczy}[ok. 2 s.]
\begin{itemize}
  \item Transformer w~zarysie.
  \item Model jako rozkład prawdopodobieństwa następnego tokenu. Dostrajanie do instrukcji.
  \item Temperatura i~ziarno: dlaczego odpowiedź jest zmienną losową i~dlaczego do pomiaru progów potrzebujemy wielu prób.
  \item Wniosek dla rozdziału 4: zachowanie mierzymy jako prawdopodobieństwo, nie pojedynczą odpowiedź.
\end{itemize}

\zrodlaplanowane{\cite{vaswani2017attention, brown2020language, ouyang2022training, renze2024effect, atil2025nondeterminism}}
\end{planroboczy}

\section{Modele rozumujące: DeepSeek-R1 i~destylacja}
\label{sec:2-2}

\begin{planroboczy}[ok. 1,5 s.]
\begin{itemize}
  \item Łańcuch rozumowania.
  \item DeepSeek-R1 trenowany przez RL.
  \item Destylaty (14B na bazie Qwen2.5).
  \item Zalecenia producenta (temperatura 0,6, brak promptu systemowego).
  \item Konsekwencje dla eksperymentu: czas wywołania, osobny kanał rozumowania.
\end{itemize}

\zrodlaplanowane{\cite{wei2022chain, guo2025deepseekr1, qwen2024qwen25, deepseek2025distillcard}}
\end{planroboczy}

\section{Agent oparty na LLM: pętla percepcja–decyzja–działanie}
\label{sec:2-3}

\begin{planroboczy}[ok. 2 s.]
\begin{itemize}
  \item Architektura agenta LLM (profil, pamięć, planowanie, działanie).
  \item ReAct, Reflexion, generatywni agenci.
  \item CoALA jako ramy pojęciowe.
  \item Umiejscowienie agenta ogrodowego: agent bez pamięci długotrwałej, decyzja w~jednym kroku.
\end{itemize}

\zrodlaplanowane{\cite{wang2024survey, sumers2024cognitive, yao2023react, shinn2023reflexion, park2023generative}}
\end{planroboczy}

\section{LLM a~motywacja: agenci z~potrzebami i~LLM jako źródło motywacji}
\label{sec:2-4}

\begin{planroboczy}[ok. 2 s.]
\begin{itemize}
  \item Dwa nurty:
  \begin{enumerate}
    \item agenci LLM z~potrzebami i~pragnieniami (Humanoid Agents, D2A) oraz zachowania „przetrwania” w~symulacjach typu Sugarscape;
    \item LLM jako źródło motywacji wewnętrznej dla agentów RL (ELLM, Motif).
  \end{enumerate}
  \item Luka: brak pomiaru progów, w~których agent LLM przełącza się na zaspokajanie potrzeby.
  \item Wyniki z~preprintów oznaczamy w~tekście jako wyniki preprintu (\texttt{03}, zasada 4): \cite{masumori2025survival}.
\end{itemize}

\zrodlaplanowane{\cite{wang2023humanoid, wang2025simulating, masumori2025survival, du2023guiding, klissarov2024motif, park2023generative}}
\end{planroboczy}

\section{Agent motywowany na bazie LLM: praca referencyjna}
\label{sec:2-5}

\begin{planroboczy}[ok. 1,5 s.]
\begin{itemize}
  \item Środowisko ogrodowe, dwa warianty promptu, wnioski Zawiślak.
  \item Rzeczowo i~bez polemiki: co pokazano i~czego nie zmierzono (brak liczbowego progu, jeden przebieg, niepowtarzalność).
  \item Szczegóły błędów trafiają do 4.1, nie tutaj.
\end{itemize}

\zrodlaplanowane{\cite{zawislak2025budowa}}
\materialyrobocze{\texttt{docs/07\_analiza\_pracy\_referencyjnej.md}}
\end{planroboczy}

    \cleardoublepage

    %% Szkielet wygenerowany z docs/02_spis_tresci.md (plan v1, 11.11.2026).
%% Blok planroboczy usuń, gdy podrozdział jest napisany.

\chapter{Wrażliwość LLM na sformułowanie promptu}
\label{ch:prompt}

% objętość planowana: ok. 5 s.

\section{Elementy promptu}
\label{sec:3-1}

\begin{planroboczy}[ok. 1 s.]
\begin{itemize}
  \item Z~czego składa się prompt agenta: rola, cel, opis stanu, lista akcji, reguły formatu.
  \item Każdy element to potencjalny czynnik eksperymentu.
\end{itemize}

\zrodlaplanowane{\cite{brown2020language, ouyang2022training, zheng2024helpful, schulhoff2024prompt, shanahan2023role}}
\end{planroboczy}

\section{Wrażliwość na format, kolejność i~drobne zmiany}
\label{sec:3-2}

\begin{planroboczy}[ok. 2 s.]
\begin{itemize}
  \item Przegląd wyników: formatowanie (duże różnice dokładności), drobne perturbacje, kolejność opcji, rozbieżności między parafrazami instrukcji, pytanie, czy modele „rozumieją” treść promptu.
  \item Wniosek metodologiczny: wynik jednego promptu nie jest miarodajny, trzeba kontrolować czynniki uboczne.
\end{itemize}

\zrodlaplanowane{\cite{sclar2024quantifying, salinas2024butterfly, pezeshkpour2024large, mizrahi2024state, webson2022prompt}}
\end{planroboczy}

\section{Ramowanie, persona i~bodźce emocjonalne}
\label{sec:3-3}

\begin{planroboczy}[ok. 1 s.]
\begin{itemize}
  \item Efekt ramowania u~ludzi.
  \item Bodźce emocjonalne w~promptach.
  \item Persona w~prompcie systemowym (brak systematycznej poprawy).
  \item Hipoteza ramowania dla progów bólu (H3).
  \item Wyniki z~preprintów oznaczamy w~tekście jako wyniki preprintu (\texttt{03}, zasada 4): \cite{li2023large}.
\end{itemize}

\zrodlaplanowane{\cite{tversky1981framing, li2023large, zheng2024helpful, shanahan2023role, sharma2024towards, jones2022capturing}}
\end{planroboczy}

\section{Kompromisy między celem a~„bólem” w~LLM}
\label{sec:3-4}

\begin{planroboczy}[ok. 1 s.]
\begin{itemize}
  \item Badanie gry punkty vs zadany ból/przyjemność: progi przełączenia u~części modeli, u~innych stałe unikanie bólu albo reakcja stopniowana.
  \item Metodologia badania LLM jak uczestnika eksperymentu psychologicznego.
  \item Najbliższy kontekst dla tej pracy.
  \item Wyniki z~preprintów oznaczamy w~tekście jako wyniki preprintu (\texttt{03}, zasada 4): \cite{keeling2024can}, \cite{bianco2026beyond}, \cite{masumori2025survival}, \cite{tagliabue2025probing}, \cite{tagliabue2026pain}, \cite{ren2026ai}, \cite{berg2026language}, \cite{schlatter2026incomplete}, \cite{mikaelson2025beyond}.
\end{itemize}

\zrodlaplanowane{\cite{keeling2024can, binz2023using, masumori2025survival, tagliabue2025probing, tagliabue2026pain, bianco2026beyond, ren2026ai, berg2026language, schlatter2026incomplete, cedro2026cash, mikaelson2025beyond}}
\end{planroboczy}

    \cleardoublepage

    %% Szkielet wygenerowany z docs/02_spis_tresci.md (plan v1, 11.11.2026).
%% Blok planroboczy usuń, gdy podrozdział jest napisany.

\chapter{Metodyka pomiaru progów bólu}
\label{ch:metodyka}

% objętość planowana: ok. 8 s.

\section{Środowisko ogrodowe i~jego modyfikacje}
\label{sec:4-1}

\begin{planroboczy}[ok. 2 s.]
\begin{itemize}
  \item Zasoby, akcje, dynamika (tabela parametrów).
  \item Kategorie słowne stanów.
  \item Odtworzenie kodu z~listingów pracy referencyjnej (jedyne źródło) i~sposób sprawdzenia rekonstrukcji (przykłady stanów z~pracy).
  \item Wprowadzone poprawki (temperatura w \texttt{options}, ziarno, lokalny generator losowy, parsowanie, logowanie, wyrównanie stan–akcja) z~uzasadnieniem.
  \item Wersja środowiska.
\end{itemize}

\zrodlaplanowane{\cite{zawislak2025budowa, starzyk2013simulation, ollama2026api, qwen2024qwen25, li2024evaluating}}
\materialyrobocze{\texttt{docs/05\_srodowisko\_i\_kod.md}}
\end{planroboczy}

\section{Operacjonalizacja: ból, akcja naprawcza, próg}
\label{sec:4-2}

\begin{planroboczy}[ok. 1,5 s.]
\begin{itemize}
  \item Definicje formalne z \texttt{04\_metodologia\_pomiaru.md}: poziom zasobu, zbiór akcji naprawczych, odpowiedź binarna, próg.
  \item Odniesienie do definicji Starzyka.
  \item Uzasadnienie podejścia „LLM jako badany”.
\end{itemize}

\zrodlaplanowane{\cite{starzyk2017needs, binz2023using, hagendorff2023machine, keeling2024can, dominguezolmedo2024questioning}}
\end{planroboczy}

\section{Funkcja psychometryczna i~estymacja}
\label{sec:4-3}

\begin{planroboczy}[ok. 2 s.]
\begin{itemize}
  \item Model z~parametrami $\theta$, $k$, $\gamma$, $\lambda$.
  \item Estymacja metodą największej wiarygodności.
  \item Bootstrap, porównanie wariantów promptu.
  \item Wzory (2–4, numerowane) i~rysunek przykładowej krzywej z~oznaczeniem parametrów.
\end{itemize}

\zrodlaplanowane{\cite{wichmann2001psychometric, wichmann2001bootstrap, kingdom2016psychophysics, miller2024adding, prins2018applying, schutt2016painfree, kuss2005bayesian, holm1979simple}}
\end{planroboczy}

\section{Protokół sondowania kontrolowanego}
\label{sec:4-4}

\begin{planroboczy}[ok. 1 s.]
\begin{itemize}
  \item Stany sztuczne: jeden zasób zmienny, pozostałe neutralne.
  \item Liczba poziomów i~prób.
  \item Temperatura i~ziarna, tryb rozumowania, parsowanie.
\end{itemize}

\zrodlaplanowane{\cite{deepseek2025distillcard, renze2024effect, atil2025nondeterminism, song2025good, ollama2026api}}
\end{planroboczy}

\section{Protokół pętli zamkniętej i~metryki zachowania}
\label{sec:4-5}

\begin{planroboczy}[ok. 1 s.]
\begin{itemize}
  \item Epizody naprawcze (początek i~koniec), histereza.
  \item Opóźnienie reakcji, udział kroków w~bólu, akcje bezskuteczne, bezczynność.
  \item Wspólne ziarna dla wariantów.
  \item Wyniki z~preprintów oznaczamy w~tekście jako wyniki preprintu (\texttt{03}, zasada 4): \cite{masumori2025survival}.
\end{itemize}

\zrodlaplanowane{\cite{keramati2014homeostatic, astrom2021feedback, park2023generative, masumori2025survival, liu2024agentbench, zawislak2025budowa}}
\materialyrobocze{\texttt{docs/04\_metodologia\_pomiaru.md}}
\end{planroboczy}

\section{Warianty promptów i~plan eksperymentów}
\label{sec:4-6}

\begin{planroboczy}[ok. 0,5 s.]
\begin{itemize}
  \item Tabela czynników i~poziomów (skrót z \texttt{06}).
  \item Kontrola czynników ubocznych (kolejność akcji, parafrazy).
  \item Budżet obliczeniowy.
\end{itemize}

\zrodlaplanowane{\cite{mizrahi2024state, pezeshkpour2024large, sclar2024quantifying, atil2025nondeterminism, miller2024adding, prins2018applying, song2025good, holm1979simple}}
\materialyrobocze{\texttt{docs/06\_plan\_eksperymentow.md}}
\end{planroboczy}

    \cleardoublepage

    %% Szkielet wygenerowany z docs/02_spis_tresci.md (plan v1, 11.11.2026).
%% Blok planroboczy usuń, gdy podrozdział jest napisany.

\chapter{Wyniki}
\label{ch:wyniki}

\begin{planroboczy}[ok. 12 s.]
Każdy podrozdział:
\begin{enumerate}
  \item pytanie badawcze,
  \item warunki (tabela albo zdanie),
  \item wynik (wykres i~tabela parametrów z~przedziałami ufności),
  \item krótka interpretacja bez odwołań do literatury (te trafiają do rozdziału 6).
\end{enumerate}

Liczby wyłącznie z \texttt{results/}.
\end{planroboczy}

\section{Weryfikacja narzędzia pomiarowego}
\label{sec:5-1}

\begin{planroboczy}[ok. 1 s.]
\begin{itemize}
  \item Powtarzalność przy stałym ziarnie.
  \item Odsetek odpowiedzi nieczytelnych, rozkład czasu wywołania.
  \item Porównanie z~trybem bez rozumowania (jeśli E0 to obejmie).
  \item Test odzyskiwania parametrów na danych syntetycznych.
\end{itemize}
\end{planroboczy}

\section{Progi bazowe: bateria i~wilgotność gleby}
\label{sec:5-2}

\begin{planroboczy}[ok. 2 s.]
\begin{itemize}
  \item PB1: krzywe psychometryczne dla promptu bazowego.
  \item Czy przełączenie jest progowe? Porównanie z~interpretacją jakościową z~pracy referencyjnej.
\end{itemize}
\end{planroboczy}

\section{Wpływ reprezentacji stanu}
\label{sec:5-3}

\begin{planroboczy}[ok. 1,5 s.]
\begin{itemize}
  \item PB4: etykiety vs liczby vs oba.
\end{itemize}
\end{planroboczy}

\section{Kalibracja zadanym progiem}
\label{sec:5-4}

\begin{planroboczy}[ok. 2 s.]
\begin{itemize}
  \item PB2: $\theta$ w~funkcji zadanego $X$. Nachylenie i~przesunięcie zależności.
\end{itemize}
\end{planroboczy}

\section{Ramowanie i~wskazówki zachowania}
\label{sec:5-5}

\begin{planroboczy}[ok. 2,5 s.]
\begin{itemize}
  \item PB3 i~PB5: przesunięcia $\theta$ i~zmiany $k$, $\gamma$, $\lambda$ względem promptu bazowego. Wykres „leśny” różnic z~przedziałami ufności.
\end{itemize}
\end{planroboczy}

\section{Rywalizacja bólów}
\label{sec:5-6}

\begin{planroboczy}[ok. 1,5 s.]
\begin{itemize}
  \item PB6: mapa prawdopodobieństwa wyboru bateria/gleba na siatce 2D. Granica decyzyjna dla 2–3 promptów.
\end{itemize}
\end{planroboczy}

\section{Walidacja w~pętli zamkniętej}
\label{sec:5-7}

\begin{planroboczy}[ok. 1,5 s.]
\begin{itemize}
  \item PB7: metryki epizodów naprawczych i~zależność od $\theta$ z~sondowania.
\end{itemize}
\end{planroboczy}

    \cleardoublepage

    %% Szkielet wygenerowany z docs/02_spis_tresci.md (plan v1, 11.11.2026).
%% Blok planroboczy usuń, gdy podrozdział jest napisany.

\chapter{Dyskusja}
\label{ch:dyskusja}

% objętość planowana: ok. 3,5 s.

\section{Progi agenta LLM a~model uczenia motywowanego}
\label{sec:6-1}

\begin{planroboczy}[ok. 1,5 s.]
\begin{itemize}
  \item Czy agent LLM zachowuje się jak agent z~progiem?
  \item Co odpowiada WTA, a~co bólom abstrakcyjnym.
  \item Modele progów odpowiedzi z~etologii (podział pracy u~owadów społecznych): ta sama sigmoidalna funkcja progowa z~parametrem stromości; progi stałe i~progi wzmacniane doświadczeniem jako analogia do kalibracji progu promptem (PB2); krytyka: same progi nie wyjaśniają całego zachowania (\cite{bonabeau1996quantitative}, \cite{theraulaz1998response}, \cite{ulrich2021response}).
\end{itemize}

\zrodlaplanowane{\cite{starzyk2017needs, zawislak2025budowa, graham2015opportunistic, keramati2014homeostatic, bonabeau1996quantitative, theraulaz1998response, ulrich2021response, yoshida2025linking}}
\end{planroboczy}

\section{Wyniki na tle literatury o~wrażliwości na prompt}
\label{sec:6-2}

\begin{planroboczy}[ok. 1 s.]
\begin{itemize}
  \item Porównanie skali efektów z~wcześniejszymi badaniami wrażliwości na prompt i~z~progami przełączenia z~badania punkty–ból.
  \item Wyniki z~preprintów oznaczamy w~tekście jako wyniki preprintu (\texttt{03}, zasada 4): \cite{keeling2024can}, \cite{masumori2025survival}, \cite{tagliabue2026pain}, \cite{schlatter2026incomplete}, \cite{mikaelson2025beyond}.
\end{itemize}

\zrodlaplanowane{\cite{sclar2024quantifying, keeling2024can, masumori2025survival, tagliabue2026pain, schlatter2026incomplete, mizrahi2024state, sharma2024towards, cedro2026cash, mikaelson2025beyond}}
\end{planroboczy}

\section{Ograniczenia i~kierunki dalszych badań}
\label{sec:6-3}

\begin{planroboczy}[ok. 1 s.]
\begin{itemize}
  \item Jeden model, uproszczone środowisko, dyskretne kategorie stanu, budżet prób.
  \item Ryzyko antropomorfizacji.
  \item Propozycje dla kolejnych prac, w~tym styk z~pracą Kingi.
  \item Zależność od preprintów: wątek z~3.4 opiera się głównie na preprintach (stan X.2026); w~tekście zaznaczamy, które twierdzenia ich dotyczą.
  \item Wyniki z~preprintów oznaczamy w~tekście jako wyniki preprintu (\texttt{03}, zasada 4): \cite{ren2026ai}.
\end{itemize}

\zrodlaplanowane{\cite{sharkey2025could, li2024evaluating, ren2026ai, shanahan2023role, hagendorff2023machine, atil2025nondeterminism, dominguezolmedo2024questioning, song2025good, butlin2026identifying, long2024taking}}
\end{planroboczy}

    \cleardoublepage

%---------------------------------------------------------------------------------
%                      PODSUMOWANIE
%---------------------------------------------------------------------------------
    \chapter*{Podsumowanie}
    \addcontentsline{toc}{chapter}{Podsumowanie}
    %% Szkielet wygenerowany z docs/02_spis_tresci.md (plan v1, 11.11.2026).
%% Blok planroboczy usuń, gdy podrozdział jest napisany.

\begin{planroboczy}[ok. 2 s.]
Odpowiedzi na PB1–PB7 w~punktach, wkład własny, perspektywy. Zgodnie z~wymogiem WMiFS: ocena wyników i~odniesienie do literatury.
\end{planroboczy}


    \cleardoublepage

    %-----------------------------------------------------------------------------------
    %                            BIBLIOGRAFIA
    %-----------------------------------------------------------------------------------
    % Plik generowany: python tools/bib2bibitem.py --bib docs/literatura.bib --main thesis/Praca_magisterska.tex --out thesis/bibliografia.tex
    \phantomsection\addcontentsline{toc}{chapter}{Bibliografia}
    % Plik wygenerowany przez tools/bib2bibitem.py -- nie edytować ręcznie.
\begin{thebibliography}{99}

\bibitem{asada2019artificial} M.~Asada, „Artificial Pain May Induce Empathy, Morality, and Ethics in the Conscious Mind of Robots”, \textit{Philosophies}, t.~4, nr~3, art.~38, 2019. Dostępne online: \url{https://doi.org/10.3390/philosophies4030038} [dostęp: 11.11.2026].

\bibitem{astrom2021feedback} K.~J.~Åström, R.~M.~Murray, \textit{Feedback Systems: An Introduction for Scientists and Engineers}, wyd.~2, Princeton University Press, 2021. Dostępne online: \url{https://fbswiki.org/wiki/index.php/Feedback_Systems:_An_Introduction_for_Scientists_and_Engineers} [dostęp: 08.10.2026].

\bibitem{atil2025nondeterminism} B.~Atıl i~in., „Non-Determinism of «Deterministic» {LLM} System Settings in Hosted Environments”, w:~\textit{Proceedings of the 5th Workshop on Evaluation and Comparison of NLP Systems (Eval4NLP 2025)}, s.~135--148, Association for Computational Linguistics, 2025. Dostępne online: \url{https://doi.org/10.18653/v1/2025.eval4nlp-1.12} [dostęp: 11.11.2026].

\bibitem{berg2026language} C.~Berg, C.~Kaiser, „Language Models Act on Hidden Valence”, preprint arXiv:2609.35591v1, 2026. Dostępne online: \url{https://arxiv.org/abs/2609.35591v1} [dostęp: 11.11.2026].

\bibitem{bianco2026beyond} F.~Bianco, D.~Shiller, „Beyond Behavioural Trade-Offs: Mechanistic Tracing of Pain-Pleasure Decisions in an {LLM}”, preprint arXiv:2602.19159v1, 2026. Dostępne online: \url{https://arxiv.org/abs/2602.19159v1} [dostęp: 11.11.2026].

\bibitem{binz2023using} M.~Binz, E.~Schulz, „Using Cognitive Psychology to Understand {GPT-3}”, \textit{Proceedings of the National Academy of Sciences}, t.~120, nr~6, art.~e2218523120, 2023. Dostępne online: \url{https://doi.org/10.1073/pnas.2218523120} [dostęp: 11.11.2026].

\bibitem{bonabeau1996quantitative} E.~Bonabeau, G.~Theraulaz, J.-L.~Deneubourg, „Quantitative Study of the Fixed Threshold Model for the Regulation of Division of Labour in Insect Societies”, \textit{Proceedings of the Royal Society of London B}, t.~263, nr~1376, s.~1565--1569, 1996. DOI: 10.1098/rspb.1996.0229.

\bibitem{brown2020language} T.~B.~Brown i~in., „Language Models are Few-Shot Learners”, w:~\textit{Advances in Neural Information Processing Systems 33 (NeurIPS 2020)}, s.~1877--1901, 2020. Dostępne online: \url{https://proceedings.neurips.cc/paper/2020/hash/1457c0d6bfcb4967418bfb8ac142f64a-Abstract.html} [dostęp: 08.10.2026].

\bibitem{butlin2026identifying} P.~Butlin i~in., „Identifying Indicators of Consciousness in {AI} Systems”, \textit{Trends in Cognitive Sciences}, t.~30, nr~6, s.~488--501, 2026. Dostępne online: \url{https://doi.org/10.1016/j.tics.2025.10.011} [dostęp: 11.11.2026].

\bibitem{cedro2026cash} M.~Cedro i~in., „Cash or Comfort? How {LLMs} Value Your Inconvenience”, \textit{Communications of the ACM}, t.~69, nr~9, s.~58--66, 2026. Dostępne online: \url{https://doi.org/10.1145/3799434} [dostęp: 08.10.2026].

\bibitem{deepseek2025distillcard} DeepSeek-AI, \textit{{DeepSeek-R1-Distill-Qwen-14B}: karta modelu}. Dostępne online: \url{https://huggingface.co/deepseek-ai/DeepSeek-R1-Distill-Qwen-14B} [dostęp: 27.09.2026].

\bibitem{dominguezolmedo2024questioning} R.~Dominguez-Olmedo, M.~Hardt, C.~Mendler-Dünner, „Questioning the Survey Responses of Large Language Models”, w:~\textit{Advances in Neural Information Processing Systems 37 (NeurIPS 2024)}, s.~45850--45878, 2024. Dostępne online: \url{https://proceedings.neurips.cc/paper_files/paper/2024/hash/515c62809e0a29729d7eec26e2916fc0-Abstract-Conference.html} [dostęp: 08.10.2026].

\bibitem{du2023guiding} Y.~Du i~in., „Guiding Pretraining in Reinforcement Learning with Large Language Models”, w:~\textit{Proceedings of the 40th International Conference on Machine Learning}, Proceedings of Machine Learning Research, t.~202, s.~8657--8677, 2023. Dostępne online: \url{https://proceedings.mlr.press/v202/du23f.html} [dostęp: 08.10.2026].

\bibitem{graham2015opportunistic} J.~Graham, J.~A.~Starzyk, D.~Jachyra, „Opportunistic Behavior in Motivated Learning Agents”, \textit{IEEE Transactions on Neural Networks and Learning Systems}, t.~26, nr~8, s.~1735--1746, 2015. DOI: 10.1109/TNNLS.2014.2354400.

\bibitem{guo2025deepseekr1} D.~Guo i~in., „{DeepSeek-R1} Incentivizes Reasoning in {LLMs} through Reinforcement Learning”, \textit{Nature}, t.~645, nr~8081, s.~633--638, 2025. Dostępne online: \url{https://doi.org/10.1038/s41586-025-09422-z} [dostęp: 11.11.2026].

\bibitem{hagendorff2023machine} T.~Hagendorff i~in., „Machine Psychology”, preprint arXiv:2303.13988v6, 2024. Dostępne online: \url{https://arxiv.org/abs/2303.13988v6} [dostęp: 11.11.2026].

\bibitem{holm1979simple} S.~Holm, „A Simple Sequentially Rejective Multiple Test Procedure”, \textit{Scandinavian Journal of Statistics}, t.~6, nr~2, s.~65--70, 1979. \url{https://www.jstor.org/stable/4615733}.

\bibitem{jones2022capturing} E.~Jones, J.~Steinhardt, „Capturing Failures of Large Language Models via Human Cognitive Biases”, w:~\textit{Advances in Neural Information Processing Systems 35 (NeurIPS 2022)}, s.~11785--11799, 2022. Dostępne online: \url{https://proceedings.neurips.cc/paper_files/paper/2022/hash/4d13b2d99519c5415661dad44ab7edcd-Abstract-Conference.html} [dostęp: 08.10.2026].

\bibitem{keeling2024can} G.~Keeling i~in., „Can {LLMs} Make Trade-offs Involving Stipulated Pain and Pleasure States?”, preprint arXiv:2411.02432v1, 2024. Dostępne online: \url{https://arxiv.org/abs/2411.02432v1} [dostęp: 11.11.2026].

\bibitem{keramati2014homeostatic} M.~Keramati, B.~Gutkin, „Homeostatic Reinforcement Learning for Integrating Reward Collection and Physiological Stability”, \textit{eLife}, t.~3, art.~e04811, 2014. Dostępne online: \url{https://doi.org/10.7554/eLife.04811} [dostęp: 11.11.2026].

\bibitem{kingdom2016psychophysics} F.~A.~A.~Kingdom, N.~Prins, \textit{Psychophysics: A Practical Introduction}, wyd.~2, Academic Press, London 2016.

\bibitem{klissarov2024motif} M.~Klissarov i~in., „Motif: Intrinsic Motivation from Artificial Intelligence Feedback”, w:~\textit{International Conference on Learning Representations (ICLR 2024)}, 2024. Dostępne online: \url{https://openreview.net/forum?id=tmBKIecDE9} [dostęp: 08.10.2026].

\bibitem{kuehn2017artificial} J.~Kuehn, S.~Haddadin, „An Artificial Robot Nervous System to Teach Robots How to Feel Pain and Reflexively React to Potentially Damaging Contacts”, \textit{IEEE Robotics and Automation Letters}, t.~2, nr~1, s.~72--79, 2017. DOI: 10.1109/LRA.2016.2536360.

\bibitem{kuss2005bayesian} M.~Kuss, F.~Jäkel, F.~A.~Wichmann, „{Bayesian} Inference for Psychometric Functions”, \textit{Journal of Vision}, t.~5, nr~5, s.~478--492, 2005. Dostępne online: \url{https://doi.org/10.1167/5.5.8} [dostęp: 11.11.2026].

\bibitem{li2023large} C.~Li i~in., „Large Language Models Understand and Can Be Enhanced by Emotional Stimuli”, preprint arXiv:2307.11760v7, 2023. Dostępne online: \url{https://arxiv.org/abs/2307.11760v7} [dostęp: 11.11.2026].

\bibitem{li2024evaluating} S.~Li i~in., „Evaluating Quantized Large Language Models”, w:~\textit{Proceedings of the 41st International Conference on Machine Learning}, Proceedings of Machine Learning Research, t.~235, s.~28480--28524, 2024. Dostępne online: \url{https://proceedings.mlr.press/v235/li24bb.html} [dostęp: 11.11.2026].

\bibitem{liu2024agentbench} X.~Liu i~in., „{AgentBench}: Evaluating {LLMs} as Agents”, w:~\textit{International Conference on Learning Representations (ICLR 2024)}, 2024. Dostępne online: \url{https://openreview.net/forum?id=zAdUB0aCTQ} [dostęp: 08.10.2026].

\bibitem{long2024taking} R.~Long i~in., „Taking {AI} Welfare Seriously”, preprint arXiv:2411.00986v1, 2024. Dostępne online: \url{https://arxiv.org/abs/2411.00986v1} [dostęp: 11.11.2026].

\bibitem{man2019homeostasis} K.~Man, A.~Damasio, „Homeostasis and Soft Robotics in the Design of Feeling Machines”, \textit{Nature Machine Intelligence}, t.~1, nr~10, s.~446--452, 2019. DOI: 10.1038/s42256-019-0103-7.

\bibitem{maslow1943theory} A.~H.~Maslow, „A Theory of Human Motivation”, \textit{Psychological Review}, t.~50, nr~4, s.~370--396, 1943. DOI: 10.1037/h0054346. Dostępne online: \url{https://psychclassics.yorku.ca/Maslow/motivation.htm} [dostęp: 08.10.2026].

\bibitem{masumori2025survival} A.~Masumori, T.~Ikegami, „Do Large Language Model Agents Exhibit a Survival Instinct? {An} Empirical Study in a {Sugarscape}-Style Simulation”, preprint arXiv:2508.12920v1, 2025. Dostępne online: \url{https://arxiv.org/abs/2508.12920v1} [dostęp: 11.11.2026].

\bibitem{mikaelson2025beyond} L.~Mikaelson, D.~Shiller, H.~Clatterbuck, „Beyond Mimicry: Preference Coherence in {LLMs}”, preprint arXiv:2511.13630v1, 2025. Dostępne online: \url{https://arxiv.org/abs/2511.13630v1} [dostęp: 11.11.2026].

\bibitem{miller2024adding} E.~Miller, „Adding Error Bars to Evals: A Statistical Approach to Language Model Evaluations”, preprint arXiv:2411.00640v1, 2024. Dostępne online: \url{https://arxiv.org/abs/2411.00640v1} [dostęp: 11.11.2026].

\bibitem{mizrahi2024state} M.~Mizrahi i~in., „State of What Art? {A} Call for Multi-Prompt {LLM} Evaluation”, \textit{Transactions of the Association for Computational Linguistics}, t.~12, s.~933--949, 2024. Dostępne online: \url{https://doi.org/10.1162/tacl_a_00681} [dostęp: 11.11.2026].

\bibitem{ollama2026api} Ollama, \textit{Ollama {API} Reference: Generate a Chat Completion}. Dostępne online: \url{https://github.com/ollama/ollama/blob/main/docs/api.md} [dostęp: 27.09.2026].

\bibitem{oudeyer2007intrinsic} P.-Y.~Oudeyer, F.~Kaplan, „What is Intrinsic Motivation? A Typology of Computational Approaches”, \textit{Frontiers in Neurorobotics}, t.~1, art.~6, 2007. Dostępne online: \url{https://doi.org/10.3389/neuro.12.006.2007} [dostęp: 11.11.2026].

\bibitem{ouyang2022training} L.~Ouyang i~in., „Training Language Models to Follow Instructions with Human Feedback”, w:~\textit{Advances in Neural Information Processing Systems 35 (NeurIPS 2022)}, s.~27730--27744, 2022. Dostępne online: \url{https://proceedings.neurips.cc/paper_files/paper/2022/hash/b1efde53be364a73914f58805a001731-Abstract-Conference.html} [dostęp: 08.10.2026].

\bibitem{park2023generative} J.~S.~Park i~in., „Generative Agents: Interactive Simulacra of Human Behavior”, w:~\textit{Proceedings of the 36th Annual ACM Symposium on User Interface Software and Technology (UIST '23)}, s.~1--22, ACM, 2023. Dostępne online: \url{https://doi.org/10.1145/3586183.3606763} [dostęp: 11.11.2026].

\bibitem{pezeshkpour2024large} P.~Pezeshkpour, E.~Hruschka, „Large Language Models Sensitivity to the Order of Options in Multiple-Choice Questions”, w:~\textit{Findings of the Association for Computational Linguistics: NAACL 2024}, s.~2006--2017, Association for Computational Linguistics, 2024. Dostępne online: \url{https://doi.org/10.18653/v1/2024.findings-naacl.130} [dostęp: 11.11.2026].

\bibitem{prins2018applying} N.~Prins, F.~A.~A.~Kingdom, „Applying the Model-Comparison Approach to Test Specific Research Hypotheses in Psychophysical Research Using the Palamedes Toolbox”, \textit{Frontiers in Psychology}, t.~9, art.~1250, 2018. Dostępne online: \url{https://doi.org/10.3389/fpsyg.2018.01250} [dostęp: 11.11.2026].

\bibitem{ren2026ai} R.~Ren i~in., „{AI} Wellbeing: Measuring and Improving the Functional Pleasure and Pain of {AIs}”, manuskrypt, Center for AI Safety, wersja PDF z 2.10.2026, 2026. Dostępne online: \url{https://www.ai-wellbeing.org/paper.pdf} [dostęp: 08.10.2026].

\bibitem{renze2024effect} M.~Renze, „The Effect of Sampling Temperature on Problem Solving in Large Language Models”, w:~\textit{Findings of the Association for Computational Linguistics: EMNLP 2024}, s.~7346--7356, Association for Computational Linguistics, 2024. Dostępne online: \url{https://doi.org/10.18653/v1/2024.findings-emnlp.432} [dostęp: 11.11.2026].

\bibitem{russell2020artificial} S.~Russell, P.~Norvig, \textit{Artificial Intelligence: A Modern Approach}, wyd.~4, Pearson, Hoboken 2020.

\bibitem{ryan2000self} R.~M.~Ryan, E.~L.~Deci, „Self-Determination Theory and the Facilitation of Intrinsic Motivation, Social Development, and Well-Being”, \textit{American Psychologist}, t.~55, nr~1, s.~68--78, 2000. DOI: 10.1037/0003-066X.55.1.68.

\bibitem{salinas2024butterfly} A.~Salinas, F.~Morstatter, „The Butterfly Effect of Altering Prompts: How Small Changes and Jailbreaks Affect Large Language Model Performance”, w:~\textit{Findings of the Association for Computational Linguistics: ACL 2024}, s.~4629--4651, Association for Computational Linguistics, 2024. Dostępne online: \url{https://doi.org/10.18653/v1/2024.findings-acl.275} [dostęp: 11.11.2026].

\bibitem{schlatter2026incomplete} J.~Schlatter, B.~Weinstein-Raun, J.~Ladish, „Incomplete Tasks Induce Shutdown Resistance in Some Frontier {LLMs}”, preprint arXiv:2509.14260v2, 2026. Dostępne online: \url{https://arxiv.org/abs/2509.14260v2} [dostęp: 11.11.2026].

\bibitem{schulhoff2024prompt} S.~Schulhoff i~in., „The Prompt Report: A Systematic Survey of Prompt Engineering Techniques”, preprint arXiv:2406.06608v6, 2024. Dostępne online: \url{https://arxiv.org/abs/2406.06608v6} [dostęp: 11.11.2026].

\bibitem{schutt2016painfree} H.~H.~Schütt i~in., „Painfree and Accurate {Bayesian} Estimation of Psychometric Functions for (Potentially) Overdispersed Data”, \textit{Vision Research}, t.~122, s.~105--123, 2016. Dostępne online: \url{https://doi.org/10.1016/j.visres.2016.02.002} [dostęp: 11.11.2026].

\bibitem{sclar2024quantifying} M.~Sclar i~in., „Quantifying Language Models' Sensitivity to Spurious Features in Prompt Design or: How {I} Learned to Start Worrying about Prompt Formatting”, w:~\textit{International Conference on Learning Representations (ICLR 2024)}, 2024. Dostępne online: \url{https://openreview.net/forum?id=RIu5lyNXjT} [dostęp: 08.10.2026].

\bibitem{shanahan2023role} M.~Shanahan, K.~McDonell, L.~Reynolds, „Role Play with Large Language Models”, \textit{Nature}, t.~623, nr~7987, s.~493--498, 2023. Dostępne online: \url{https://doi.org/10.1038/s41586-023-06647-8} [dostęp: 11.11.2026].

\bibitem{sharkey2025could} A.~Sharkey, „Could a Robot Feel Pain?”, \textit{AI \& Society}, t.~40, nr~5, s.~3641--3651, 2025. Dostępne online: \url{https://doi.org/10.1007/s00146-024-02110-y} [dostęp: 11.11.2026].

\bibitem{sharma2024towards} M.~Sharma i~in., „Towards Understanding Sycophancy in Language Models”, w:~\textit{International Conference on Learning Representations (ICLR 2024)}, 2024. Dostępne online: \url{https://openreview.net/forum?id=tvhaxkMKAn} [dostęp: 08.10.2026].

\bibitem{shinn2023reflexion} N.~Shinn i~in., „Reflexion: Language Agents with Verbal Reinforcement Learning”, w:~\textit{Advances in Neural Information Processing Systems 36 (NeurIPS 2023)}, s.~8634--8652, 2023. Dostępne online: \url{https://proceedings.neurips.cc/paper_files/paper/2023/hash/1b44b878bb782e6954cd888628510e90-Abstract-Conference.html} [dostęp: 08.10.2026].

\bibitem{singh2010intrinsically} S.~Singh i~in., „Intrinsically Motivated Reinforcement Learning: An Evolutionary Perspective”, \textit{IEEE Transactions on Autonomous Mental Development}, t.~2, nr~2, s.~70--82, 2010. DOI: 10.1109/TAMD.2010.2051031. Dostępne online (wersja autorska): \url{https://web.eecs.umich.edu/~baveja/Papers/IMRLIEEETAMDFinal.pdf} [dostęp: 08.10.2026].

\bibitem{song2025good} Y.~Song i~in., „The Good, The Bad, and The Greedy: Evaluation of {LLMs} Should Not Ignore Non-Determinism”, w:~\textit{Proceedings of NAACL 2025 (Long Papers)}, s.~4195--4206, 2025. Dostępne online: \url{https://doi.org/10.18653/v1/2025.naacl-long.211} [dostęp: 11.11.2026].

\bibitem{starzyk2008motivation} J.~A.~Starzyk, „Motivation in Embodied Intelligence”, w:~\textit{Frontiers in Robotics, Automation and Control}, I-Tech Education and Publishing, 2008, s.~83--110. Dostępne online: \url{https://doi.org/10.5772/6332} [dostęp: 11.11.2026].

\bibitem{starzyk2011motivated} J.~A.~Starzyk, „Motivated Learning for Computational Intelligence”, w:~B.~Igelnik (red.), \textit{Computational Modeling and Simulation of Intellect: Current State and Future Perspectives}, IGI Publishing, Hershey 2011, s.~265--292. DOI: 10.4018/978-1-60960-551-3.ch011. Dostępne online (wersja autorska): \url{http://ace.cs.ohio.edu/~starzyk/network/Research/Papers/Motivated%20Learning%20for%20Computational%20Intelligence.pdf} [dostęp: 08.10.2026].

\bibitem{starzyk2012motivated} J.~A.~Starzyk i~in., „Motivated Learning for the Development of Autonomous Systems”, \textit{Cognitive Systems Research}, t.~14, nr~1, s.~10--25, 2012. DOI: 10.1016/j.cogsys.2010.12.009. Dostępne online (preprint): \url{http://ace.cs.ohio.edu/~starzyk/network/Research/Papers/Motivated_Learning_for_Autonomous_Development.pdf} [dostęp: 08.10.2026].

\bibitem{starzyk2013simulation} J.~A.~Starzyk, J.~Graham, L.~Puzio, „Simulation of a Motivated Learning Agent”, w:~\textit{Artificial Intelligence Applications and Innovations (AIAI 2013)}, IFIP Advances in Information and Communication Technology, t.~412, s.~205--214, Springer, 2013. Dostępne online: \url{https://doi.org/10.1007/978-3-642-41142-7_21} [dostęp: 11.11.2026].

\bibitem{starzyk2017mlecog} J.~A.~Starzyk, J.~Graham, „{MLECOG}: Motivated Learning Embodied Cognitive Architecture”, \textit{IEEE Systems Journal}, t.~11, nr~3, s.~1272--1283, 2017. DOI: 10.1109/JSYST.2015.2442995.

\bibitem{starzyk2017needs} J.~A.~Starzyk, J.~Graham, L.~Puzio, „Needs, Pains, and Motivations in Autonomous Agents”, \textit{IEEE Transactions on Neural Networks and Learning Systems}, t.~28, nr~11, s.~2528--2540, 2017. DOI: 10.1109/TNNLS.2016.2596787.

\bibitem{sumers2024cognitive} T.~R.~Sumers i~in., „Cognitive Architectures for Language Agents”, \textit{Transactions on Machine Learning Research}, 2024. Dostępne online: \url{https://arxiv.org/abs/2309.02427} [dostęp: 11.11.2026].

\bibitem{sutton2018reinforcement} R.~S.~Sutton, A.~G.~Barto, \textit{Reinforcement Learning: An Introduction}, wyd.~2, MIT Press, Cambridge, MA 2018. Dostępne online: \url{http://incompleteideas.net/book/RLbook2020.pdf} [dostęp: 08.10.2026].

\bibitem{tagliabue2025probing} V.~Tagliabue, L.~Dung, „Probing the Preferences of a Language Model: Integrating Verbal and Behavioral Tests of {AI} Welfare”, preprint arXiv:2509.07961v2, 2025. Dostępne online: \url{https://arxiv.org/abs/2509.07961v2} [dostęp: 11.11.2026].

\bibitem{tagliabue2026pain} V.~Tagliabue, L.~Dung, C.~Berg, „The Pain Axis: {LLMs} Represent Self-Directed Harm and Act on It”, preprint arXiv:2609.16247v2, 2026. Dostępne online: \url{https://arxiv.org/abs/2609.16247v2} [dostęp: 11.11.2026].

\bibitem{theraulaz1998response} G.~Theraulaz, E.~Bonabeau, J.-L.~Deneubourg, „Response Threshold Reinforcements and Division of Labour in Insect Societies”, \textit{Proceedings of the Royal Society of London B}, t.~265, nr~1393, s.~327--332, 1998. DOI: 10.1098/rspb.1998.0299. Dostępne online: \url{https://pmc.ncbi.nlm.nih.gov/articles/PMC1688885/} [dostęp: 08.10.2026].

\bibitem{tversky1981framing} A.~Tversky, D.~Kahneman, „The Framing of Decisions and the Psychology of Choice”, \textit{Science}, t.~211, nr~4481, s.~453--458, 1981. DOI: 10.1126/science.7455683.

\bibitem{ulrich2021response} Y.~Ulrich i~in., „Response Thresholds Alone Cannot Explain Empirical Patterns of Division of Labor in Social Insects”, \textit{PLOS Biology}, t.~19, nr~6, art.~e3001269, 2021. Dostępne online: \url{https://doi.org/10.1371/journal.pbio.3001269} [dostęp: 11.11.2026].

\bibitem{vaswani2017attention} A.~Vaswani i~in., „Attention Is All You Need”, w:~\textit{Advances in Neural Information Processing Systems 30 (NeurIPS 2017)}, s.~5998--6008, 2017. Dostępne online: \url{https://proceedings.neurips.cc/paper_files/paper/2017/hash/3f5ee243547dee91fbd053c1c4a845aa-Abstract.html} [dostęp: 08.10.2026].

\bibitem{wang2023humanoid} Z.~Wang, Y.~Y.~Chiu, Y.~C.~Chiu, „Humanoid Agents: Platform for Simulating Human-like Generative Agents”, w:~\textit{Proceedings of the 2023 Conference on Empirical Methods in Natural Language Processing: System Demonstrations}, s.~167--176, Association for Computational Linguistics, 2023. Dostępne online: \url{https://doi.org/10.18653/v1/2023.emnlp-demo.15} [dostęp: 11.11.2026].

\bibitem{wang2024survey} L.~Wang i~in., „A Survey on Large Language Model Based Autonomous Agents”, \textit{Frontiers of Computer Science}, t.~18, nr~6, art.~186345, 2024. Dostępne online: \url{https://doi.org/10.1007/s11704-024-40231-1} [dostęp: 11.11.2026].

\bibitem{wang2025simulating} Y.~Wang i~in., „Simulating Human-like Daily Activities with Desire-driven Autonomy”, w:~\textit{International Conference on Learning Representations (ICLR 2025)}, 2025. Dostępne online: \url{https://openreview.net/forum?id=3ms8EQY7f8} [dostęp: 08.10.2026].

\bibitem{webson2022prompt} A.~Webson, E.~Pavlick, „Do Prompt-Based Models Really Understand the Meaning of Their Prompts?”, w:~\textit{Proceedings of the 2022 Conference of the North American Chapter of the Association for Computational Linguistics: Human Language Technologies}, s.~2300--2344, Association for Computational Linguistics, 2022. Dostępne online: \url{https://doi.org/10.18653/v1/2022.naacl-main.167} [dostęp: 11.11.2026].

\bibitem{wei2022chain} J.~Wei i~in., „Chain-of-Thought Prompting Elicits Reasoning in Large Language Models”, w:~\textit{Advances in Neural Information Processing Systems 35 (NeurIPS 2022)}, s.~24824--24837, 2022. Dostępne online: \url{https://proceedings.neurips.cc/paper_files/paper/2022/hash/9d5609613524ecf4f15af0f7b31abca4-Abstract-Conference.html} [dostęp: 08.10.2026].

\bibitem{wichmann2001psychometric} F.~A.~Wichmann, N.~J.~Hill, „The Psychometric Function: {I}. {Fitting}, Sampling, and Goodness of Fit”, \textit{Perception \& Psychophysics}, t.~63, nr~8, s.~1293--1313, 2001. Dostępne online: \url{https://doi.org/10.3758/BF03194544} [dostęp: 11.11.2026].

\bibitem{wichmann2001bootstrap} F.~A.~Wichmann, N.~J.~Hill, „The Psychometric Function: {II}. {Bootstrap}-Based Confidence Intervals and Sampling”, \textit{Perception \& Psychophysics}, t.~63, nr~8, s.~1314--1329, 2001. Dostępne online: \url{https://doi.org/10.3758/BF03194545} [dostęp: 11.11.2026].

\bibitem{qwen2024qwen25} A.~Yang i~in., „{Qwen2.5} Technical Report”, preprint arXiv:2412.15115v2, 2024. Dostępne online: \url{https://arxiv.org/abs/2412.15115v2} [dostęp: 11.11.2026].

\bibitem{yao2023react} S.~Yao i~in., „{ReAct}: Synergizing Reasoning and Acting in Language Models”, w:~\textit{International Conference on Learning Representations (ICLR 2023)}, 2023. Dostępne online: \url{https://openreview.net/forum?id=WE_vluYUL-X} [dostęp: 08.10.2026].

\bibitem{yoshida2025linking} N.~Yoshida, H.~Sprekeler, B.~Gutkin, „Linking homeostasis to reinforcement learning: internal state control of motivated behavior”, \textit{Current Opinion in Behavioral Sciences}, t.~66, art.~101611, 2025. Dostępne online: \url{https://doi.org/10.1016/j.cobeha.2025.101611} [dostęp: 11.11.2026].

\bibitem{zawislak2025budowa} K.~Zawiślak, \textit{Budowa agenta motywowanego na bazie Dużego Modelu Językowego ({LLM})}, praca magisterska, Politechnika Rzeszowska im. Ignacego Łukasiewicza, Wydział Matematyki i Fizyki Stosowanej, Rzeszów 2025.

\bibitem{zheng2024helpful} M.~Zheng i~in., „When «{A} Helpful Assistant» Is Not Really Helpful: Personas in System Prompts Do Not Improve Performances of Large Language Models”, w:~\textit{Findings of the Association for Computational Linguistics: EMNLP 2024}, s.~15126--15154, Association for Computational Linguistics, 2024. Dostępne online: \url{https://doi.org/10.18653/v1/2024.findings-emnlp.888} [dostęp: 11.11.2026].

\end{thebibliography}

    \cleardoublepage

    %-----------------------------------------------------------------------------------
    %                            SPISY RYSUNKÓW/TABEL
    %-----------------------------------------------------------------------------------
    %\renewcommand{\listfigurename}{\LARGE Spis rysunków}
    %\addcontentsline{toc}{chapter}{Spis rysunków}
    %\listoffigures
    %\cleardoublepage

    %\renewcommand{\listtablename}{\LARGE Spis tabel}
    %\addcontentsline{toc}{chapter}{Spis tabel}
    %\listoftables
    %\cleardoublepage

    %-----------------------------------------------------------------------------------
    %                            DODATKI/ANEKSY
    %-----------------------------------------------------------------------------------
	

    %% Szkielet wygenerowany z docs/02_spis_tresci.md (plan v1, 11.11.2026).
%% Blok planroboczy usuń, gdy podrozdział jest napisany.

\chapter*{Dodatek A. Pełne treści promptów}
\addcontentsline{toc}{chapter}{Dodatek A. Pełne treści promptów}
\label{app:a}

\begin{planroboczy}
Prompty wstawiane z~plików \texttt{code/prompts/} (np. \texttt{\textbackslash lstinputlisting}), żeby praca pokazywała dokładnie to, co dostał model.
\end{planroboczy}

\cleardoublepage
\chapter*{Dodatek B. Parametry środowiska i~kategorie stanów}
\addcontentsline{toc}{chapter}{Dodatek B. Parametry środowiska i~kategorie stanów}
\label{app:b}

\begin{planroboczy}
Tabele parametrów środowiska ogrodowego i~granic kategorii stanów.
\end{planroboczy}

\cleardoublepage
\chapter*{Dodatek C. Instrukcja odtworzenia wyników}
\addcontentsline{toc}{chapter}{Dodatek C. Instrukcja odtworzenia wyników}
\label{app:c}

\begin{planroboczy}
Repozytorium, wersje oprogramowania i~modelu, polecenia.
\end{planroboczy}

\cleardoublepage
\chapter*{Dodatek D. Wyniki uzupełniające}
\addcontentsline{toc}{chapter}{Dodatek D. Wyniki uzupełniające}
\label{app:d}

\begin{planroboczy}
Krzywe psychometryczne dla wszystkich wariantów promptu.
\end{planroboczy}

    

    %-----------------------------------------------------------------------------------
    %                            STRESZCZENIA
    %-----------------------------------------------------------------------------------
	

    \cleardoublepage
	\thispagestyle{empty}
\setlength{\parindent}{0in}
\setlength{\parskip}{0em}

\textbf{POLITECHNIKA RZESZOWSKA} \hfill Rzeszów 2027 \\
\textbf{im. Ignacego Łukasiewicza} \\
\textbf{WYDZIAŁ MATEMATYKI I FIZYKI STOSOWANEJ} \\

\begin{center}
	\textbf{STRESZCZENIE  PRACY  DYPLOMOWEJ }
\end{center}

\textbf{Tytuł:} Wpływ promptu na progi bólu agenta LLM \\
\textbf{Autor:} Kamil Piejko \\
\textbf{Promotor:} dr inż. Marcin Kowalik, prof. PRz \\
\textbf{Słowa klucze:} \dower{do 5 słów kluczowych}
\\

\dower{tekst streszczenia w języku polskim, pisany na końcu}


% Wstawienie pionowej przestrzeni do połowy strony
\vfill
$$\noindent\rule{\textwidth}{2pt}$$



\textbf{RZESZOW UNIVERSITY OF TECHNOLOGY} \hfill Rzeszów 2027 \\
\textbf{FACULTY OF MATHEMATICS AND APPLIED PHYSICS } \\

\begin{center}
	\textbf{DIPLOMA THESIS ABSTRACT}
\end{center}

\textbf{Title:} The Influence of Prompt Wording on the Pain Thresholds of an LLM Agent \dower{tytuł angielski: propozycja, do zatwierdzenia}\\
\textbf{Author:} Kamil Piejko \\
\textbf{Supervisor:} Marcin Kowalik, PhD Eng., Assoc. Prof. \dower{zapis stopnia promotora po angielsku do potwierdzenia} \\ %% Angielski zapis stopni/tytułów naukowych promotora 

\textbf{Key words:} \dower{up to 5 key words}
\\

\dower{tekst streszczenia w języku angielskim}

\vfill

% Poniżej przykłady angielskich odpowiedników polskich stopni i tytułów naukowych promotora:
% - doktor | dr → Dr Firstname Lastname lub Firstname Lastname, PhD
% - dr inż. → Dr Eng. Firstname Lastname lub Firstname Lastname, PhD Eng.
% - doktor habilitowany | dr hab. → Firstname Lastname, PhD, DSc
% - dr hab. inż. → Firstname Lastname, PhD, DSc Eng.
% - dr hab. prof. PRz → Firstname Lastname, PhD, DSc, Assoc. Prof.
% - prof. dr hab. → Prof. Firstname Lastname
% - prof. dr hab. inż. → Prof. Eng. Firstname Lastname
% - doktor, profesor PRz → Doctor, Firstname Lastname, Associate Professor lub Firstname Lastname, PhD, Assoc. Prof.

    \thispagestyle{empty}

    %-----------------------------------------------------------------------------------
    %                            TABELA GENAI
    %-----------------------------------------------------------------------------------
    \cleardoublepage
	
\thispagestyle{empty}
\begin{center}
    Wykaz obszarów i narzędzi GenAI \\
    wykorzystanych przez autora w trakcie wykonywania pracy dyplomowej
\end{center}

\begin{table}[ht]
    \centering
    \begin{tabularx}{\linewidth}{|c|X|X|}\hline
       \textbf{Lp.}  & \multicolumn{1}{c|}{\textbf{Obszar wykorzystywania $^{a)}$}} & \multicolumn{1}{c|}{\textbf{Narzędzie}} \\ \hline
         &  &\\ \hline
         &  &\\ \hline
    \end{tabularx}
    \label{tab:my_label}
\end{table}

\textbf{Uwaga!} W przypadku, gdy autor nie korzysta z obszarów i narzędzi GenAI w tabeli wpisuje „nie korzystano”

\

\textit{$^{a)}$ Przykładowe obszary wykorzystania:
    \begin{enumerate}[label=\alph*)]
        \item redakcja i korekta tekstu (m.in. poprawa stylistyczna i gramatyczna, synonimy i~parafrazowanie, sprawdzenie spójności językowej),
        \item analiza stanu wiedzy, przegląd literatury,
        \item generowanie treści lub przykładów,
        \item tworzenie pytań do ankiet i kwestionariuszy,
        \item tworzenie schematów, diagramów i map myśli,
        \item analiza danych i wizualizacja wyników (m.in. obliczenia statystyczne, tworzenie wykresów i tabel),
        \item tworzenie podsumowań i wniosków,
        \item pisanie i testowanie kodu,
        \item tłumaczenia językowe.
    \end{enumerate}} %obowiązkowe dla studentów wybierających tematy od semestru letniego 2024/2025, pozostali mogą usunąć
	
	\end{document}

	
